\documentclass[11pt]{article}

\usepackage{amsmath}
\usepackage{amssymb}
\usepackage{amsthm}
\usepackage{geometry}
\geometry{margin=1in}

\newtheorem{lemma}{Lemma}
\newtheorem{claim}{Claim}
\theoremstyle{remark}
\newtheorem*{remark}{Remark}

\title{Why the Discounted-Reward Approximation is a Lower Bound}
\author{}
\date{}

\begin{document}
\maketitle

\section*{Setup}

Let $\gamma \in (0,1)$ be the discount factor, and let $L$ denote the (random) length,
in environment timesteps, of a won episode. Let $p$ denote the pursuer's win rate.
The true expected discounted reward is
\[
  \mathbb{E}\big[\mathbb{1}\{\text{won}\}\cdot \gamma^{L}\big]
  \;=\; p \cdot \mathbb{E}\big[\gamma^{L} \,\big|\, \text{won}\big],
\]
while the approximation used when only $\texttt{ep\_len\_mean} := \mathbb{E}[L \mid \text{won}]$
is available (rather than the full distribution of $L$) is
\[
  p \cdot \gamma^{\,\mathbb{E}[L\mid\text{won}]}.
\]
Since the win-rate factor $p$ is common to both expressions, the question reduces
entirely to the conditional distribution of $L$ given a win. We show
\begin{equation}
  \mathbb{E}[\gamma^{L}] \;\ge\; \gamma^{\,\mathbb{E}[L]} \tag{$\star$}
\end{equation}
(expectations below implicitly conditioned on winning), so the approximation never
overstates the true discounted reward.

\section*{Step 1: $f(x) = \gamma^x$ is strictly convex}

\begin{lemma}
For $\gamma \in (0,1)$, the function $f(x) = \gamma^{x} = e^{x\ln\gamma}$ is strictly
convex on all of $\mathbb{R}$.
\end{lemma}

\begin{proof}
Differentiating twice,
\[
  f'(x) = (\ln\gamma)\,\gamma^{x}, \qquad
  f''(x) = (\ln\gamma)^{2}\,\gamma^{x}.
\]
Since $\gamma \in (0,1)$ we have $\ln\gamma \neq 0$, so $(\ln\gamma)^2 > 0$; and
$\gamma^{x} > 0$ for every real $x$. Hence $f''(x) > 0$ for all $x \in \mathbb{R}$,
which implies $f$ is strictly convex on $\mathbb{R}$.
\end{proof}

\section*{Step 2: Convex functions lie above their tangent lines}

\begin{lemma}[Supporting tangent line]
Let $\phi$ be convex and differentiable on $\mathbb{R}$, and fix $x_0 \in \mathbb{R}$.
Then for every $x \in \mathbb{R}$,
\[
  \phi(x) \;\ge\; \phi(x_0) + \phi'(x_0)(x - x_0).
\]
\end{lemma}

\begin{proof}
By convexity, for $t \in (0,1]$,
\[
  \phi\big((1-t)x_0 + t x\big) \;\le\; (1-t)\,\phi(x_0) + t\,\phi(x).
\]
Subtracting $\phi(x_0)$ from both sides and dividing by $t$,
\[
  \frac{\phi\big(x_0 + t(x - x_0)\big) - \phi(x_0)}{t} \;\le\; \phi(x) - \phi(x_0).
\]
The left-hand side is the difference quotient of $\phi$ at $x_0$ in the direction
$(x - x_0)$. Letting $t \to 0^{+}$ and using differentiability of $\phi$ at $x_0$,
the left-hand side converges to $\phi'(x_0)(x - x_0)$, giving
\[
  \phi'(x_0)(x - x_0) \;\le\; \phi(x) - \phi(x_0),
\]
which rearranges to the claimed inequality.
\end{proof}

\section*{Step 3: Jensen's inequality}

\begin{claim}[Jensen's inequality for $f = \gamma^{(\cdot)}$]
For any random variable $L$ with finite mean, $\mathbb{E}[\gamma^{L}] \ge \gamma^{\mathbb{E}[L]}$.
\end{claim}

\begin{proof}
Apply Step 2 with $\phi = f$, $x_0 = \mathbb{E}[L]$, and $x = L$ (for every realization
of $L$):
\[
  f(L) \;\ge\; f(\mathbb{E}[L]) + f'(\mathbb{E}[L])\big(L - \mathbb{E}[L]\big).
\]
Taking expectations of both sides and using linearity of expectation,
\[
  \mathbb{E}[f(L)] \;\ge\; f(\mathbb{E}[L]) + f'(\mathbb{E}[L])\underbrace{\big(\mathbb{E}[L] - \mathbb{E}[L]\big)}_{=\,0}
  \;=\; f(\mathbb{E}[L]).
\]
That is, $\mathbb{E}[\gamma^{L}] \ge \gamma^{\mathbb{E}[L]}$, which is $(\star)$.
Multiplying both sides by $p \ge 0$ preserves the inequality and yields the original
claim:
\[
  \mathbb{E}\big[\mathbb{1}\{\text{won}\}\cdot\gamma^{L}\big] \;\ge\; p\cdot\gamma^{\,\mathbb{E}[L\mid\text{won}]}. \qedhere
\]
\end{proof}

\section*{Equality condition}

Equality in Step 2 holds only where $\phi$ is affine between $x_0$ and $x$; since
$f'' > 0$ everywhere, $f$ is \emph{strictly} convex, so its tangent line at $x_0$
touches $f$ only at $x = x_0$ itself. Consequently, equality in $(\star)$ holds if
and only if $L = \mathbb{E}[L]$ almost surely --- every winning episode takes exactly
the same number of steps. For a stochastic policy in a stochastic environment this
never holds exactly, so the bound is strict in practice: the approximation
understates the true discounted reward, by an amount that grows with the variance
of win-episode lengths.

\begin{remark}[Two-point sanity check]
Suppose half of all wins take $a$ steps and half take $b$ steps, $a \neq b$. Then
\[
  \mathbb{E}[\gamma^{L}] = \tfrac12\gamma^{a} + \tfrac12\gamma^{b},
  \qquad
  \gamma^{\mathbb{E}[L]} = \gamma^{(a+b)/2} = \sqrt{\gamma^{a}\gamma^{b}}.
\]
This is exactly the AM--GM inequality applied to the positive numbers $\gamma^a$ and
$\gamma^b$: the arithmetic mean is always at least the geometric mean, with equality
iff $a = b$ --- matching the general equality condition above via an elementary
route.
\end{remark}

\begin{remark}[Magnitude in practice]
Since $f''(x) = (\ln\gamma)^2\gamma^x \to 0$ as $\gamma \to 1$, the convexity of
$f$ --- and hence the size of the gap in $(\star)$ --- is weak at $\gamma = 0.99$.
Combined with the short, low-variance win-episode lengths observed in this
environment (mean around 6--8 steps), this is why the approximation tracked the
exact values so closely wherever both were available: the bound is directionally
guaranteed, but its magnitude is small under these conditions.
\end{remark}

\end{document}
